\section{fmm3d\_\-mpi/tt1.cpp File Reference}
\label{tt1_8cpp}\index{fmm3d_mpi/tt1.cpp@{fmm3d\_\-mpi/tt1.cpp}}


{\tt \#include \char`\"{}fmm3d\_\-mpi.hpp\char`\"{}}\par
\subsection*{Functions}
\begin{CompactItemize}
\item 
int {\bf main} (int argc, char $\ast$$\ast$argv)
\end{CompactItemize}


\subsection{Function Documentation}
\index{tt1.cpp@{tt1.cpp}!main@{main}}
\index{main@{main}!tt1.cpp@{tt1.cpp}}
\subsubsection{\setlength{\rightskip}{0pt plus 5cm}int main (int {\em argc}, char $\ast$$\ast$ {\em argv})}\label{tt1_8cpp_30c4c19ff69d905e8e4c750b3af23485_30c4c19ff69d905e8e4c750b3af23485}




Definition at line 22 of file tt1.cpp.

References FMM3d\_\-MPI::ctr(), FMM3d\_\-MPI::evaluate(), Knl\-Mat3d\_\-MPI::knl(), p\-A, p\-C, FMM3d\_\-MPI::root\-Level(), FMM3d\_\-MPI::setup(), Kernel3d\_\-MPI::src\-DOF(), Knl\-Mat3d\_\-MPI::src\-Nor(), Knl\-Mat3d\_\-MPI::src\-Pos(), Kernel3d\_\-MPI::trg\-DOF(), and Knl\-Mat3d\_\-MPI::trg\-Pos().

\footnotesize\begin{verbatim}23 {
24   PetscInitialize(&argc,&argv,"options_0",NULL);  //PetscTruth flg = PETSC_FALSE;
25 
26   MPI_Comm comm; comm = PETSC_COMM_WORLD;
27   int mpirank;  MPI_Comm_rank(MPI_COMM_WORLD, &mpirank);
28   int mpisize;  MPI_Comm_size(MPI_COMM_WORLD, &mpisize);
29   int dim = 3;
30   srand48( (long)( time(NULL)+mpirank ) );
31   
32   PetscTruth flg = PETSC_FALSE;
33   pC( PetscPrintf(MPI_COMM_WORLD, "mpisize %d\n", mpisize) );
34   
35   //1. random src and trg data
36   int numsrc;  pC( PetscOptionsGetInt(PETSC_NULL, "-numsrc", &numsrc, &flg) ); pA(flg==PETSC_TRUE);
37   int numtrg;  pC( PetscOptionsGetInt(PETSC_NULL, "-numtrg", &numtrg, &flg) ); pA(flg==PETSC_TRUE);
38   int kt;  pC( PetscOptionsGetInt(PETSC_NULL, "-kt", &kt, &flg) ); pA(flg==PETSC_TRUE);
39   
40   vector<double> tmp(2);         tmp[0] = 1;     tmp[1] = 0.25; //coefs in the kernel, work for all examples
41   Kernel3d_MPI knl(kt, tmp);
42   
43   int lclnumsrc = (numsrc+mpirank)/mpisize;
44   int lclnumtrg = (numtrg+mpirank)/mpisize;
45   
46   //pick center
47   double R = 0.85;
48   double r = 1-R;
49   double lclctr[3];
50   for(int d=0; d<dim; d++)       lclctr[d] = R * (2.0*drand48()-1.0);
51   
52   //generate src and trg data
53   Vec srcPos;  pC( VecCreateMPI(comm, lclnumsrc*dim, PETSC_DETERMINE, &srcPos) );
54   double* srcPosarr; pC( VecGetArray(srcPos, &srcPosarr) );
55   for(int k=0; k<lclnumsrc; k++)
56          for(int d=0; d<dim; d++)
57                 srcPosarr[d+dim*k] = lclctr[d] + r*(2.0*drand48()-1.0);
58   pC( VecRestoreArray(srcPos, &srcPosarr) );
59   
60   Vec trgPos;  pC( VecCreateMPI(comm, lclnumtrg*dim, PETSC_DETERMINE, &trgPos) );
61   double* trgPosarr; pC( VecGetArray(trgPos, &trgPosarr) );
62   for(int k=0; k<lclnumtrg; k++)
63          for(int d=0; d<dim; d++)
64                 trgPosarr[d+dim*k] = lclctr[d] + r*(2.0*drand48()-1.0);
65   pC( VecRestoreArray(trgPos, &trgPosarr) );
66   
67   int srcDOF = knl.srcDOF();
68   int trgDOF = knl.trgDOF();
69   Vec srcDen;  pC( VecCreateMPI(comm, lclnumsrc*srcDOF, PETSC_DETERMINE, &srcDen) );
70   double* srcDenarr; pC( VecGetArray(srcDen, &srcDenarr) );
71   for(int k=0; k<lclnumsrc*srcDOF; k++)
72          srcDenarr[k] = drand48();
73   pC( VecRestoreArray(srcDen, &srcDenarr) );
74   
75   Vec trgVal;  pC( VecCreateMPI(comm, lclnumtrg*trgDOF, PETSC_DETERMINE, &trgVal) );
76   double zero = 0;  pC( VecSet(trgVal, zero) );
77   
78   //2. allocate fmm 
79   clock_t ck0, ck1;
80   
81   FMM3d_MPI* fmm = new FMM3d_MPI("fmm3d_MPI_");
82   fmm->srcPos()=srcPos;  fmm->srcNor()=srcPos;
83   fmm->trgPos()=trgPos;
84   fmm->ctr() = Point3(0,0,0); // CENTER OF THE TOPLEVEL BOX
85   fmm->rootLevel() = 0;         // 2^(-rootlvl) is the RADIUS OF THE TOPLEVEL BOX
86   fmm->knl() = knl;
87   //setup
88   ck0 = clock();
89   pC( fmm->setup() );
90   ck1 = clock();  
91   pC( PetscPrintf(MPI_COMM_WORLD, "fmm setup used %e secs\n", double(ck1-ck0)/CLOCKS_PER_SEC) );
92   
93   //3. run fmm
94   for(int i=0; i<3; i++) {
95          ck0 = clock();
96          pC( fmm->evaluate(srcDen, trgVal) );
97          //pC( MPI_Barrier(MPI_COMM_WORLD) );  //check Vincnt, Votcnt
98          ck1 = clock();
99          pC( PetscPrintf(MPI_COMM_WORLD, "fmm evaluate used %e secs\n", double(ck1-ck0)/CLOCKS_PER_SEC) );
100   }
101   
102   //4. check
103   ck0 = clock();
104   double rerr;
105   pC( fmm->check(srcDen, trgVal, 20, rerr) );
106   pC( PetscPrintf(MPI_COMM_WORLD, "relative %e\n", rerr) );
107   ck1 = clock();
108   pC( PetscPrintf(MPI_COMM_WORLD, "fmm check used %e secs\n", double(ck1-ck0)/CLOCKS_PER_SEC) );
109   //extra
110   //{ PetscViewer viewer;  pC( PetscViewerASCIIOpen(comm, "srcPos.out", &viewer) );  pC( VecView(srcPos, viewer) );  pC( PetscViewerDestroy(viewer) ); }
111   //{ PetscViewer viewer;  pC( PetscViewerASCIIOpen(comm, "trgPos.out", &viewer) );  pC( VecView(trgPos, viewer) );  pC( PetscViewerDestroy(viewer) ); }
112   
113   delete fmm;
114   
115   pC( VecDestroy(srcPos) );
116   pC( VecDestroy(trgPos) );
117   pC( VecDestroy(srcDen) );
118   pC( VecDestroy(trgVal) );
119   
120   PetscFinalize();
121   return 0;
122 }
\end{verbatim}\normalsize 


